% Setting up the document and importing necessary packages
\documentclass[12pt,a4paper]{article}

% Chinese support packages
\usepackage{xeCJK}
\usepackage{fontspec}
\setCJKmainfont{Noto Serif CJK TC}
\setCJKsansfont{Noto Sans CJK TC}
\setCJKmonofont{Noto Sans CJK TC}

% Other necessary packages
\usepackage{geometry}
\usepackage{titlesec}
\usepackage{enumitem}
\usepackage{color}
\usepackage{colortbl}
\usepackage{tcolorbox}
\usepackage{booktabs}
\usepackage{longtable}
\usepackage{array}
\usepackage{graphicx}
\usepackage{tikz}
\usepackage{mdframed}
\usepackage{fancyhdr}
\usepackage{hyperref}
\usepackage{multicol}
\usepackage{datetime2}
\usepackage{natbib}

\hypersetup{
    colorlinks=true,
    linkcolor=heme_red,
    citecolor=heme_red,
    urlcolor=heme_red,
    pdfborder={0 0 0}
}

% Page setup
\geometry{
    top=2.5cm,
    bottom=2.5cm,
    left=2cm,
    right=2cm,
    headheight=15pt
}

% Color definitions (Hematology themed colors)
\definecolor{heme_red}{RGB}{178,34,52}
\definecolor{heme_crimson}{RGB}{220,20,60}
\definecolor{heme_gray}{RGB}{120,120,120}
\definecolor{highlight_yellow}{RGB}{255,250,205}

% Title format settings
\titleformat{\section}[block]{\Large\bfseries\color{heme_red}}{\thesection}{10pt}{}
\titleformat{\subsection}[block]{\large\bfseries\color{heme_crimson}}{\thesubsection}{8pt}{}
\titleformat{\subsubsection}[block]{\normalsize\bfseries\color{heme_gray}}{\thesubsubsection}{6pt}{}

% Custom environment
\newmdenv[
    linecolor=heme_red,
    backgroundcolor=highlight_yellow,
    linewidth=2pt,
    topline=true,
    bottomline=true,
    leftline=true,
    rightline=true,
    innertopmargin=10pt,
    innerbottommargin=10pt,
    innerleftmargin=10pt,
    innerrightmargin=10pt
]{importantbox}

\newcommand{\note}[1]{
    \par
    \noindent
    \begin{tcolorbox}[
        colback=heme_crimson!5!white,
        colframe=heme_crimson,
        title=\textbf{重點提醒},
        fonttitle=\bfseries,
        boxrule=0.5pt,
        arc=3pt,
        left=6pt,
        right=6pt,
        top=6pt,
        bottom=6pt
    ]
    #1
    \end{tcolorbox}
}

% Header and footer settings
\pagestyle{fancy}
\fancyhf{}
\fancyhead[L]{\textcolor{heme_red}{內科部住院醫師教學文件}}
\fancyhead[R]{\textcolor{heme_crimson}{血液科EPAs}}
\fancyfoot[C]{\textcolor{heme_gray}{\thepage}}
\renewcommand{\headrulewidth}{1pt}
\renewcommand{\headrule}{\hbox to\headwidth{\color{heme_red}\leaders\hrule height \headrulewidth\hfill}}

% Date format
\DTMsetstyle{yyyy-mm-dd}
\newcommand{\mydate}{\DTMtoday}

% Document start
\begin{document}

% Title page
\begin{titlepage}
    \centering
    {\LARGE\textcolor{heme_red}{\textbf{內科部住院醫師教學文件}}\par}
    \vspace{1cm}
    {\huge\bfseries 血液科可信賴專業活動EPAs\par}
    \vspace{0.5cm}
    {\Large\textcolor{heme_crimson}{\textit{Hematology Entrustable Professional Activities}}\par}
    \vspace{1cm}
    
    \begin{tcolorbox}[
        colback=heme_crimson!5!white,
        colframe=heme_crimson,
        width=0.8\textwidth,
        arc=10pt,
        boxrule=1pt
    ]
    \centering
    {\large\textbf{目標對象}\par}
    \vspace{0.3cm}
    內科部住院醫師R1-R3\par
    血液科專科訓練醫師\par
    \vspace{0.5cm}
    {\large\textbf{文件版本}\par}
    \vspace{0.3cm}
    第1.0版(10個EPAs)\par
    發布日期:\mydate\par
    \end{tcolorbox}

    \vfill
    
    {\large 血液科 \\
    適用於CCC評量系統\par}
    
    \vspace{0.5cm}
\end{titlepage}

\newpage
\tableofcontents
\newpage

% Main content
\section{血液科EPAs概述}

\subsection{EPA定義}
Entrustable Professional Activities (EPA) 是指可信賴委託給住院醫師執行的專業活動單元。血液科EPAs涵蓋了血液疾病診斷、治療與照護的核心能力,是所有血液科醫師必須精熟的專業技能。

\vspace{0.5cm}
\note{
血液科EPAs評量著重於能否安全、穩定地執行完整的血液疾病診療流程,並能整合資訊做出適當臨床決策。每個EPA都包含明確的里程碑發展階段,從初學者到專家的完整能力建構。本版本包含10個核心EPA,從基礎的病史詢問與血液抹片判讀,到複雜的骨髓檢查、化療管理與跨科部協調,提供完整的血液疾病診療能力架構。
}

\subsection{學習目標}
完成血液科EPAs後,住院醫師應能夠:
\begin{itemize}[nosep]
    \item 熟練進行血液疾病病史詢問與理學檢查。
    \item 判讀周邊血液抹片與骨髓抹片。
    \item 診斷與處理各類型貧血。
    \item 診斷與管理出血性疾病與凝血異常。
    \item 診斷與處理血栓性疾病。
    \item 執行骨髓穿刺與切片檢查。
    \item 診斷與治療急性白血病。
    \item 診斷與管理淋巴瘤與骨髓瘤。
    \item 執行輸血醫學與成分輸血。
    \item 協調血液疾病患者跨科部整合照護。
\end{itemize}

\subsection{委託等級}
\begin{table}[h]
    \centering
    \caption{EPA委託等級定義}
    \begin{tabular}{p{2cm} p{3cm} p{9cm}}
        \toprule
        \textbf{等級} & \textbf{委託程度} & \textbf{描述} \\
        \midrule
        Level 1 & 觀察學習 & 僅能觀察他人執行,無法獨立操作 \\
        Level 2 & 直接監督 & 可在主治醫師直接監督下執行 \\
        Level 3 & 間接監督 & 可獨立執行,但需回報並接受指導 \\
        Level 4 & 監督可得 & 可獨立執行,監督者隨時可得 \\
        Level 5 & 完全獨立 & 可獨立執行並指導他人 \\
        \bottomrule
    \end{tabular}
\end{table}

\newpage
\section{血液科10大EPAs詳述}

\begin{longtable}{|p{1.2cm}|p{3.8cm}|p{8.5cm}|}
\caption{血液科EPAs完整架構} \\
\hline
\rowcolor{heme_crimson!20}
\textbf{EPA} & \textbf{專業活動名稱} & \textbf{核心能力要素與里程碑} \\
\hline
\endfirsthead

\multicolumn{3}{c}%
{{\bfseries \tablename\ \thetable{} -- 接續上頁}} \\
\hline
\rowcolor{heme_crimson!20}
\textbf{EPA} & \textbf{專業活動名稱} & \textbf{核心能力要素與里程碑} \\
\hline
\endhead

\hline \multicolumn{3}{|r|}{{接續下頁}} \\ \hline
\endfoot

\hline \hline
\endlastfoot

\rowcolor{heme_red!10}
\multicolumn{3}{|c|}{\textbf{EPA 1: 血液疾病病史詢問與理學檢查}} \\
\hline
\textbf{描述} & \multicolumn{2}{|p{12.5cm}|}{能夠獨立進行血液疾病相關的完整病史詢問與理學檢查,包括出血傾向、貧血症狀、淋巴結腫大評估,並能初步評估疾病嚴重度。} \\
\hline
\textbf{核心能力} & 出血症狀評估 & 瘀青、紫斑、牙齦出血、鼻出血、月經過多詳細詢問 \\
\cline{2-3}
& 貧血症狀 & 疲倦、暈眩、心悸、呼吸困難、運動耐力評估 \\
\cline{2-3}
& 感染史 & 反覆感染、發燒型態、嚴重度評估 \\
\cline{2-3}
& 家族史 & 血友病、地中海貧血、血栓傾向家族史 \\
\cline{2-3}
& 理學檢查 & 蒼白、黃疸、淋巴結、肝脾腫大、瘀斑檢查 \\
\cline{2-3}
& 功能評估 & ECOG score、Karnofsky score評估 \\
\hline
\textbf{Level 1} & Novice & 能詢問基本血液症狀,需要監導者引導完成 \\
\hline
\textbf{Level 2} & Advanced Beginner & 能系統性完成血液病史詢問,執行基本理學檢查 \\
\hline
\textbf{Level 3} & Competent & 能獨立完成複雜病史詢問,整合症狀進行初步診斷 \\
\hline
\textbf{Level 4} & Proficient & 能在困難情況下獲取準確病史,識別罕見症狀 \\
\hline
\textbf{Level 5} & Expert & 能指導他人進行病史詢問,建立系統性評估流程 \\
\hline

\rowcolor{heme_red!10}
\multicolumn{3}{|c|}{\textbf{EPA 2: 周邊血液抹片與骨髓抹片判讀}} \\
\hline
\textbf{描述} & \multicolumn{2}{|p{12.5cm}|}{能夠獨立判讀周邊血液抹片與骨髓抹片,識別正常與異常血球型態,並能整合臨床資訊提出診斷建議。} \\
\hline
\textbf{核心能力} & 紅血球型態 & 大小不等、形狀異常、內含物識別 \\
\cline{2-3}
& 白血球分類 & 五分類計數、異常白血球識別 \\
\cline{2-3}
& 血小板評估 & 數量、大小、型態評估 \\
\cline{2-3}
& 骨髓細胞 & M:E ratio、成熟度、異常細胞識別 \\
\cline{2-3}
& 特殊染色 & 鐵染色、PAS、MPO等判讀 \\
\cline{2-3}
& 臨床關聯 & 抹片發現與疾病診斷關聯性 \\
\hline
\textbf{Level 1} & Novice & 能識別基本正常血球型態 \\
\hline
\textbf{Level 2} & Advanced Beginner & 能識別常見異常血球,進行基本分類 \\
\hline
\textbf{Level 3} & Competent & 能獨立判讀大部分抹片,整合臨床診斷 \\
\hline
\textbf{Level 4} & Proficient & 能判讀複雜抹片,識別罕見血球異常 \\
\hline
\textbf{Level 5} & Expert & 成為抹片判讀專家,能教導複雜判讀技巧 \\
\hline

\rowcolor{heme_red!10}
\multicolumn{3}{|c|}{\textbf{EPA 3: 貧血診斷與處理}} \\
\hline
\textbf{描述} & \multicolumn{2}{|p{12.5cm}|}{能夠診斷各類型貧血,執行系統性診斷評估,制定個別化治療計畫,並監測治療效果與副作用。} \\
\hline
\textbf{核心能力} & 分類診斷 & MCV分類、網狀球計數、系統性診斷流程 \\
\cline{2-3}
& 缺鐵性貧血 & 診斷標準、鐵劑治療、原因追查 \\
\cline{2-3}
& 巨球性貧血 & B12/葉酸缺乏、診斷與治療 \\
\cline{2-3}
& 溶血性貧血 & 內因性vs外因性、實驗室診斷 \\
\cline{2-3}
& 慢性病貧血 & 發炎、腎病、內分泌相關貧血 \\
\cline{2-3}
& 輸血指標 & 輸血適應症、目標血紅素評估 \\
\hline
\textbf{Level 1} & Novice & 能識別貧血,知道基本分類 \\
\hline
\textbf{Level 2} & Advanced Beginner & 能診斷常見貧血,執行基本治療 \\
\hline
\textbf{Level 3} & Competent & 能獨立診斷與治療各類貧血 \\
\hline
\textbf{Level 4} & Proficient & 能處理複雜貧血,整合多重病因 \\
\hline
\textbf{Level 5} & Expert & 成為貧血診斷專家,處理罕見病因 \\
\hline

\rowcolor{heme_red!10}
\multicolumn{3}{|c|}{\textbf{EPA 4: 出血性疾病與凝血異常診斷處理}} \\
\hline
\textbf{描述} & \multicolumn{2}{|p{12.5cm}|}{能夠診斷出血性疾病與凝血異常,執行凝血功能檢查判讀,制定止血治療策略,並處理急性出血事件。} \\
\hline
\textbf{核心能力} & 凝血檢查判讀 & PT/aPTT/TT解讀、mixing study應用 \\
\cline{2-3}
& 血小板疾病 & ITP、TTP、HUS診斷與治療 \\
\cline{2-3}
& 凝血因子缺乏 & 血友病、VWD診斷與治療 \\
\cline{2-3}
& DIC診斷處理 & 診斷標準、嚴重度評估、治療策略 \\
\cline{2-3}
& 止血治療 & 血小板輸注、FFP、凝血因子使用 \\
\cline{2-3}
& 急性出血處置 & 出血評估、緊急止血、輸血決策 \\
\hline
\textbf{Level 1} & Novice & 能識別出血症狀,了解基本檢查 \\
\hline
\textbf{Level 2} & Advanced Beginner & 能判讀凝血檢查,處理簡單出血 \\
\hline
\textbf{Level 3} & Competent & 能獨立診斷凝血異常,制定治療計畫 \\
\hline
\textbf{Level 4} & Proficient & 能處理複雜出血事件,整合多重因素 \\
\hline
\textbf{Level 5} & Expert & 成為凝血專家,處理罕見凝血疾病 \\
\hline

\rowcolor{heme_red!10}
\multicolumn{3}{|c|}{\textbf{EPA 5: 血栓性疾病診斷與抗凝治療}} \\
\hline
\textbf{描述} & \multicolumn{2}{|p{12.5cm}|}{能夠診斷血栓性疾病,評估血栓風險,選擇適當抗凝藥物,監測治療效果,並處理抗凝相關併發症。} \\
\hline
\textbf{核心能力} & DVT/PE診斷 & 臨床評估、D-dimer、影像診斷 \\
\cline{2-3}
& 血栓傾向篩檢 & 遺傳性vs後天性、檢查時機 \\
\cline{2-3}
& 抗凝藥物選擇 & Warfarin、DOAC、LMWH適應症 \\
\cline{2-3}
& INR監測 & 目標範圍、劑量調整、交互作用 \\
\cline{2-3}
& 出血風險評估 & HAS-BLED score、預防策略 \\
\cline{2-3}
& 抗凝併發症 & 出血處理、HIT識別與處理 \\
\hline
\textbf{Level 1} & Novice & 能識別血栓症狀,了解基本抗凝藥物 \\
\hline
\textbf{Level 2} & Advanced Beginner & 能診斷DVT/PE,開始抗凝治療 \\
\hline
\textbf{Level 3} & Competent & 能獨立管理抗凝治療,監測與調整 \\
\hline
\textbf{Level 4} & Proficient & 能處理複雜血栓,評估血栓傾向 \\
\hline
\textbf{Level 5} & Expert & 成為血栓專家,制定個人化抗凝策略 \\
\hline

\rowcolor{heme_red!10}
\multicolumn{3}{|c|}{\textbf{EPA 6: 骨髓穿刺與切片檢查}} \\
\hline
\textbf{描述} & \multicolumn{2}{|p{12.5cm}|}{能夠執行骨髓穿刺與切片檢查,包括適應症評估、解剖定位、無菌技術、標本處理,並能初步判讀結果。} \\
\hline
\textbf{核心能力} & 適應症評估 & 診斷vs分期vs追蹤、禁忌症識別 \\
\cline{2-3}
& 解剖定位 & 後上髂骨棘、胸骨穿刺定位 \\
\cline{2-3}
& 穿刺技術 & 局部麻醉、穿刺深度、抽吸技巧 \\
\cline{2-3}
& 切片技術 & Jamshidi針使用、標本完整性 \\
\cline{2-3}
& 標本處理 & 抹片製作、固定液使用、特殊檢查 \\
\cline{2-3}
& 併發症處理 & 出血、疼痛、感染預防與處理 \\
\hline
\textbf{Level 1} & Novice & 能協助骨髓檢查,了解基本解剖 \\
\hline
\textbf{Level 2} & Advanced Beginner & 能在監督下執行骨髓穿刺 \\
\hline
\textbf{Level 3} & Competent & 能獨立執行骨髓檢查,處理標本 \\
\hline
\textbf{Level 4} & Proficient & 能執行困難穿刺,使用進階技術 \\
\hline
\textbf{Level 5} & Expert & 成為骨髓檢查專家,能教導複雜技術 \\
\hline

\rowcolor{heme_red!10}
\multicolumn{3}{|c|}{\textbf{EPA 7: 急性白血病診斷與治療}} \\
\hline
\textbf{描述} & \multicolumn{2}{|p{12.5cm}|}{能夠診斷急性白血病,執行分類與風險分層,制定化療計畫,處理化療併發症,並監測治療反應。} \\
\hline
\textbf{核心能力} & 診斷與分類 & AML vs ALL、FAB/WHO分類、免疫分型 \\
\cline{2-3}
& 細胞遺傳學 & 染色體核型、分子標記、預後評估 \\
\cline{2-3}
& 風險分層 & 預後因子、治療強度決策 \\
\cline{2-3}
& 誘導化療 & 化療方案選擇、劑量調整 \\
\cline{2-3}
& 併發症處理 & TLS、感染、出血、mucositis管理 \\
\cline{2-3}
& 反應評估 & CR標準、MRD監測、復發處理 \\
\hline
\textbf{Level 1} & Novice & 能識別白血病可能性,了解基本分類 \\
\hline
\textbf{Level 2} & Advanced Beginner & 能執行初步評估,協助化療管理 \\
\hline
\textbf{Level 3} & Competent & 能獨立診斷白血病,執行化療計畫 \\
\hline
\textbf{Level 4} & Proficient & 能處理複雜併發症,調整治療策略 \\
\hline
\textbf{Level 5} & Expert & 成為白血病專家,發展個人化治療 \\
\hline

\rowcolor{heme_red!10}
\multicolumn{3}{|c|}{\textbf{EPA 8: 淋巴瘤與骨髓瘤診斷管理}} \\
\hline
\textbf{描述} & \multicolumn{2}{|p{12.5cm}|}{能夠診斷淋巴瘤與多發性骨髓瘤,執行分期評估,制定治療計畫,監測治療反應與疾病進展。} \\
\hline
\textbf{核心能力} & 淋巴瘤診斷 & HL vs NHL、病理分型、免疫組化 \\
\cline{2-3}
& 分期評估 & Ann Arbor分期、PET-CT應用、IPI score \\
\cline{2-3}
& 化療方案 & R-CHOP、ABVD等方案選擇與執行 \\
\cline{2-3}
& 骨髓瘤診斷 & CRAB準則、ISS分期、細胞遺傳學 \\
\cline{2-3}
& 骨髓瘤治療 & 誘導、鞏固、維持治療策略 \\
\cline{2-3}
& 反應評估 & Lugano criteria、IMWG criteria應用 \\
\hline
\textbf{Level 1} & Novice & 能識別淋巴瘤可能性,了解基本分型 \\
\hline
\textbf{Level 2} & Advanced Beginner & 能執行初步分期,協助化療管理 \\
\hline
\textbf{Level 3} & Competent & 能獨立診斷與治療淋巴瘤、骨髓瘤 \\
\hline
\textbf{Level 4} & Proficient & 能處理複雜病例,評估移植適應症 \\
\hline
\textbf{Level 5} & Expert & 成為淋巴瘤專家,參與臨床試驗 \\
\hline

\rowcolor{heme_red!10}
\multicolumn{3}{|c|}{\textbf{EPA 9: 輸血醫學與成分輸血}} \\
\hline
\textbf{描述} & \multicolumn{2}{|p{12.5cm}|}{能夠評估輸血適應症,選擇適當血品,執行輸血前檢查,監測輸血反應,並處理輸血相關併發症。} \\
\hline
\textbf{核心能力} & 輸血適應症 & 紅血球、血小板、血漿輸注指標 \\
\cline{2-3}
& ABO/Rh系統 & 血型判定、交叉試驗、不規則抗體 \\
\cline{2-3}
& 成分輸血 & PRBC、PC、FFP、Cryo適應症 \\
\cline{2-3}
& 特殊輸血 & 大量輸血、緊急輸血、換血術 \\
\cline{2-3}
& 輸血反應 & 急性溶血、TRALI、TACO識別處理 \\
\cline{2-3}
& 輸血替代 & 自體輸血、無血管理策略 \\
\hline
\textbf{Level 1} & Novice & 能了解輸血基本原則,執行輸血醫囑 \\
\hline
\textbf{Level 2} & Advanced Beginner & 能評估輸血適應症,處理常規輸血 \\
\hline
\textbf{Level 3} & Competent & 能獨立管理輸血,識別輸血反應 \\
\hline
\textbf{Level 4} & Proficient & 能處理大量輸血,管理複雜輸血反應 \\
\hline
\textbf{Level 5} & Expert & 成為輸血專家,制定輸血策略 \\
\hline

\rowcolor{heme_red!10}
\multicolumn{3}{|c|}{\textbf{EPA 10: 血液疾病跨科部整合照護}} \\
\hline
\textbf{描述} & \multicolumn{2}{|p{12.5cm}|}{能夠協調血液疾病患者的整體照護,包括跨科部會診、造血幹細胞移植評估、緩和照護,並能有效溝通與教育病患及家屬。} \\
\hline
\textbf{核心能力} & 會診協調 & 感染科、放射腫瘤科、外科會診安排 \\
\cline{2-3}
& 移植評估 & 自體vs異體移植適應症評估 \\
\cline{2-3}
& 支持療法 & G-CSF、抗生素預防、營養支持 \\
\cline{2-3}
& 緩和照護 & 症狀控制、疼痛管理、善終照護 \\
\cline{2-3}
& 病患教育 & 疾病衛教、化療副作用、自我照護 \\
\cline{2-3}
& 家屬溝通 & 病情解釋、治療選項、預後討論 \\
\cline{2-3}
& 倫理考量 & 知情同意、DNR討論、醫療決策 \\
\hline
\textbf{Level 1} & Novice & 能參與基本病患照護,了解團隊角色 \\
\hline
\textbf{Level 2} & Advanced Beginner & 能執行基本會診,參與團隊討論 \\
\hline
\textbf{Level 3} & Competent & 能獨立協調病患照護,制定整合計畫 \\
\hline
\textbf{Level 4} & Proficient & 能整合複雜多科照護,領導團隊討論 \\
\hline
\textbf{Level 5} & Expert & 成為照護協調專家,改善照護品質 \\
\hline

\end{longtable}

\newpage
\section{評量方法與標準}

\subsection{CCC評量架構}
本血液科EPAs採用Case-based Clinical Competency (CCC)評量方式:
\begin{itemize}[nosep]
    \item 真實臨床情境評量
    \item 多次觀察與回饋
    \item 漸進式委託決策
    \item 360度回饋機制
\end{itemize}

\begin{importantbox}
\textbf{特別提醒:} 血液科EPAs的評量重點在於漸進式能力發展,而非單次完美表現。鼓勵從錯誤中學習,持續改進。每個EPA都需要在真實臨床情境中反復練習與評量。血液抹片判讀與骨髓檢查需要大量練習才能精熟。
\end{importantbox}

\subsection{評量工具對應表}
\begin{table}[h]
    \centering
    \caption{各EPA推薦評量工具}
    \begin{tabular}{p{4.5cm} p{4cm} p{5.5cm}}
        \toprule
        \textbf{EPA類型} & \textbf{主要評量工具} & \textbf{輔助評量方法} \\
        \midrule
        技能操作類 & DOPS & 直接觀察、技能檢核表 \\
        (EPA 6) & (Direct Observation of Procedural Skills) & 同儕評量、自我評量 \\
        \midrule
        臨床推理類 & Mini-CEX & 病例討論、口試 \\
        (EPA 1, 3, 4, 5, 7, 8, 9) & (Mini-Clinical Evaluation Exercise) & 病歷撰寫評量、模擬情境 \\
        \midrule
        抹片判讀類 & 抹片判讀評量 & 線上測驗、定期考核 \\
        (EPA 2) & Slide Interpretation Assessment & 盲測、peer review \\
        \midrule
        整合照護類 & Multi-source Feedback & 跨科團隊評量 \\
        (EPA 10) & (360-degree Assessment) & 病人滿意度調查 \\
        \bottomrule
    \end{tabular}
\end{table}

\subsection{評量頻率建議}
\begin{itemize}[nosep]
    \item \textbf{每月評量}:選擇3-4個EPA重點評量
    \item \textbf{季度檢討}:全面檢視所有10個EPA進展
    \item \textbf{年度認證}:EPA達標認證與進階規劃
    \item \textbf{即時回饋}:發現學習需求時立即介入
\end{itemize}

\section{實施建議與學習資源}

\subsection{月度晨會Mini-OSCE整合建議}
\begin{table}[h]
    \centering
    \caption{月度EPA評量主題安排}
    \begin{tabular}{p{2.5cm} p{4.5cm} p{7cm}}
        \toprule
        \textbf{月份} & \textbf{重點EPAs} & \textbf{評量重點} \\
        \midrule
        第1個月 & EPA 1, 2 & 基礎評估與抹片判讀 \\
        第2個月 & EPA 3, 9 & 貧血診斷與輸血醫學 \\
        第3個月 & EPA 4, 5 & 凝血異常與血栓管理 \\
        第4個月 & EPA 6 & 骨髓檢查技能 \\
        第5個月 & EPA 7, 8 & 血液惡性腫瘤管理 \\
        第6個月 & EPA 10 & 整合照護與協調能力 \\
        第7-12個月 & 綜合評量與進階訓練 & 複雜個案整合處理能力 \\
        \bottomrule
    \end{tabular}
\end{table}

\subsection{推薦學習資源}
學員可參考以下文獻深入學習:

\textbf{基礎教科書:}
\begin{itemize}
    \item Williams Hematology. 10th edition.
    \item Wintrobe's Clinical Hematology. 14th edition.
    \item Hoffman's Hematology: Basic Principles and Practice. 8th edition.
\end{itemize}

\textbf{臨床指引:}
\begin{itemize}
    \item ASH Guidelines for Management of Venous Thromboembolism
    \item NCCN Guidelines for Acute Lymphoblastic Leukemia
    \item NCCN Guidelines for Multiple Myeloma
    \item BSH Guidelines for Blood Transfusion
    \item EHA Guidelines for Diagnosis and Treatment of Lymphomas
\end{itemize}

\textbf{線上資源:}
\begin{itemize}
    \item ASH Image Bank (血液抹片圖庫)
    \item Hematology.org Educational Resources
    \item UpToDate Hematology Section
    \item Blood Journal Online
\end{itemize}

\subsection{自我評估工具}
\begin{itemize}[nosep]
    \item 定期記錄學習歷程與反思
    \item 建立個人抹片判讀資料庫
    \item 參與同儕學習與討論
    \item 利用ASH Image Bank練習
    \item 建立個人學習檔案
    \item 參與多專科團隊會議
    \item 定期複習骨髓檢查技能
    \item 參加血液科個案討論會
\end{itemize}

\section{總結}

血液科EPAs涵蓋了血液疾病診療的完整核心能力,從基礎的病史詢問與血液抹片判讀,到複雜的白血病化療管理、骨髓移植評估,以及整合性的跨科部協調照護。每個EPA都有清楚的能力描述與里程碑發展路徑,協助住院醫師循序漸進地建立專業能力。

\begin{importantbox}
\textbf{關鍵訊息:} 血液科EPAs的核心在於「可委託性」與「完整性」。這10個EPA代表血液科醫師必須具備的基礎診療能力。血液抹片與骨髓抹片判讀是血液科的基石,需要大量練習累積經驗。血液惡性腫瘤的治療需要謹慎評估風險效益,密切監測併發症。輸血醫學與凝血管理攸關病患安全,必須精確掌握適應症與併發症處理。
\end{importantbox}

% References section
\section{參考文獻}
\begin{thebibliography}{10}

\bibitem{williams2021}
Kaushansky, K., et al. (2021). \textit{Williams Hematology}. 10th edition. McGraw-Hill Education.

\bibitem{wintrobe2018}
Greer, J. P., et al. (2018). \textit{Wintrobe's Clinical Hematology}. 14th edition. Lippincott Williams \& Wilkins.

\bibitem{olle2013}
ten Cate, O. (2013). Nuts and bolts of entrustable professional activities. \textit{Journal of Graduate Medical Education}, 5(1), 157-158.

\bibitem{ash2020vte}
American Society of Hematology. (2020). ASH Guidelines on Management of Venous Thromboembolism. \textit{Blood Advances}, 4(19), 4693-4738.

\bibitem{nccn2024}
NCCN Clinical Practice Guidelines in Oncology. (2024). Acute Lymphoblastic Leukemia. Version 1.2024.

\bibitem{kumar2016}
Kumar, S. K., et al. (2016). Multiple myeloma. \textit{Nature Reviews Disease Primers}, 2, 16046.

\bibitem{swerdlow2016}
Swerdlow, S. H., et al. (2016). The 2016 revision of the World Health Organization classification of lymphoid neoplasms. \textit{Blood}, 127(20), 2375-2390.

\end{thebibliography}

\end{document}